% Generated draft: requires mathematical and source review.
\documentclass[11pt]{article}
\usepackage[margin=25mm]{geometry}
\usepackage{fontspec}
\setmainfont{Latin Modern Roman}
\usepackage{amsmath,amssymb,mathtools,unicode-math}
\setmathfont{Latin Modern Math}
\usepackage{enumitem,needspace,xurl,hyperref}
\setlength{\parindent}{0pt}
\setlength{\parskip}{5pt}
\setlength{\emergencystretch}{3em}
\newcommand{\N}{\mathbb{N}}
\newcommand{\Z}{\mathbb{Z}}
\newcommand{\Q}{\mathbb{Q}}
\newcommand{\R}{\mathbb{R}}
\newcommand{\C}{\mathbb{C}}
\newcommand{\F}{\mathbb{F}}
\newcommand{\A}{\mathbb{A}}
\newcommand{\PP}{\mathbb{P}}
\newcommand{\E}{\mathbb{E}}
\newcommand{\eps}{\varepsilon}
\newcommand{\dd}{\,\mathrm{d}}
\begin{document}
\section*{UnsolvedMath ID: 800234}
\textbf{Source problem code:} catalogue entry
\subsection{Definitions and mathematical statement}
Let \(\Omega\subset\C^n\) be a bounded domain, meaning a bounded, nonempty, connected open set. Say that \(\Omega\) is pseudoconvex if it admits a continuous plurisubharmonic exhaustion function. No smoothness assumption on \(\partial\Omega\) is imposed.

For \(1<p<2\) and \(0\le r\le n\), let \(L^p_{(0,r)}(\Omega)\) denote the space of differential forms of type \((0,r)\) whose coefficients belong to \(L^p(\Omega)\), with respect to Lebesgue volume. The operator \(\bar\partial\) is understood in the sense of distributions.

For a fixed degree \(1\le q\le n\), a standard meaning of \(L^p\) solvability is the assertion
\[
\forall f\in L^p_{(0,q)}(\Omega),\qquad
\bar\partial f=0
\ \Longrightarrow\ 
\exists u\in L^p_{(0,q-1)}(\Omega)
\text{ such that }\bar\partial u=f.
\]
The source asks to solve \(\bar\partial\) in \(L^p\) on bounded pseudoconvex domains, without necessarily smooth boundary, throughout \(1<p<2\).

\textit{Unresolved scope.} The source does not specify the complex dimension, the form degrees, or whether the requested conclusion includes an estimate
\[
\|u\|_{L^p(\Omega)}\le C\|f\|_{L^p(\Omega)}
\]
with \(C\) independent of \(f\), or a bounded linear solution operator. The displayed solvability assertion is therefore an explicit interpretation for fixed \(n,q,p,\Omega\); the intended quantification over dimensions and degrees and the required quantitative conclusion need review. No boundary condition on \(u\) is supplied.
\subsection{Short English statement}
For \(1<p<2\), seek \(L^p\) solutions of the \(\bar\partial\) equation on bounded pseudoconvex domains whose boundaries need not be smooth. Under the standard distributional interpretation, every \(\bar\partial\)-closed \(L^p\) datum should have an \(L^p\) primitive. The intended dimensions, form degrees, and estimates are unspecified.
\subsection{Sources}
\begin{itemize}
\item Local UnsolvedMath registry, permanent ID \(800234\), source code \texttt{catalogue entry}, titled \emph{Several Complex Variables Problem List --- Problem 31}. The supplied record labels the problem open; this is a supplied status, not an independently verified assessment.
\item Buzzard, Fornaess, and Krantz, \emph{The Several Complex Variables Problem List} (1995), Problem 31. Supplied source URL: \url{https://arxiv.org/abs/math/9509202}. The registry attributes the problem to Buzzard, Fornaess, and Krantz, with proposed year 1995.
\item Supplied provenance: accessed 2026-10-16; extraction method \texttt{arxiv\_latex}. These are supplied provenance statements; the source was not independently visited or verified for this draft.
\end{itemize}
\end{document}
